% =========================================================
% 02_vector_operations.tex
% Vectors and Linear Algebra
% =========================================================

% =========================================================
% SECTION DIVIDER
% =========================================================

\section{Dot Products and Geometry}

\begin{frame}{Dot Product}

For two vectors
\[
\vect{u}
=
\begin{bmatrix}
u_1\\
u_2\\
\vdots\\
u_n
\end{bmatrix},
\qquad
\vect{v}
=
\begin{bmatrix}
v_1\\
v_2\\
\vdots\\
v_n
\end{bmatrix},
\]

their \textbf{dot product} is
\[
\boxed{
\vect{u}\cdot\vect{v}
=
\sum_{i=1}^{n}u_i v_i
}
\]

or equivalently,

\[
\boxed{
\vect{u}\cdot\vect{v}
=
\vect{u}^{T}\vect{v}
}
\]

\end{frame}


% =========================================================
% EXAMPLE
% =========================================================

\begin{frame}{Dot Product: Example}

Consider

\[
\vect{u}
=
\begin{bmatrix}
2\\
3
\end{bmatrix},
\qquad
\vect{v}
=
\begin{bmatrix}
4\\
-1
\end{bmatrix}.
\]

Then

\[
\begin{aligned}
\vect{u}\cdot\vect{v}
&=
(2)(4)+(3)(-1)\\
&=8-3\\
&=5.
\end{aligned}
\]

Therefore,

\[
\boxed{\vect{u}\cdot\vect{v}=5}
\]

The dot product of two vectors is a \textbf{scalar}, not a vector.

\end{frame}


% =========================================================
% GEOMETRIC FORM
% =========================================================

\begin{frame}{Geometric Interpretation of the Dot Product}

The dot product can also be expressed as

\[
\boxed{
\vect{u}\cdot\vect{v}
=
\|\vect{u}\|
\|\vect{v}\|
\cos\theta
}
\]

where $\theta$ is the angle between the vectors.

\vspace{0.4cm}

Therefore,

\[
\boxed{
\cos\theta
=
\frac{
\vect{u}\cdot\vect{v}
}{
\|\vect{u}\|\|\vect{v}\|
}
}
\]

\vspace{0.3cm}

The dot product therefore connects:

\[
\text{algebra}
\quad\longleftrightarrow\quad
\text{geometry}.
\]

\end{frame}


% =========================================================
% ANGLE BETWEEN VECTORS
% =========================================================

\begin{frame}{Angle Between Two Vectors}

For non-zero vectors $\vect{u}$ and $\vect{v}$,

\[
\boxed{
\theta
=
\cos^{-1}
\left(
\frac{
\vect{u}\cdot\vect{v}
}{
\|\vect{u}\|\|\vect{v}\|
}
\right)
}
\]

\vspace{0.5cm}

The sign of the dot product provides useful geometric information:

\begin{itemize}
    \item $\vect{u}\cdot\vect{v}>0$: acute angle
    \item $\vect{u}\cdot\vect{v}=0$: right angle
    \item $\vect{u}\cdot\vect{v}<0$: obtuse angle
\end{itemize}

\end{frame}


% =========================================================
% ORTHOGONALITY
% =========================================================

\begin{frame}{Orthogonal Vectors}

Two vectors are \textbf{orthogonal} if they are perpendicular.

For vectors $\vect{u}$ and $\vect{v}$,

\vspace{-0.5cm}
\[
\boxed{
\vect{u}\cdot\vect{v}=0
}
\]

implies

\vspace{-0.5cm}
\[
\boxed{
\vect{u}\perp\vect{v}
}
\]

provided both vectors are non-zero.

\begin{exampleblock}{Example}
\[
\vect{u}
=
\begin{bmatrix}
2\\
1
\end{bmatrix},
\qquad
\vect{v}
=
\begin{bmatrix}
1\\
-2
\end{bmatrix}
\]

\vspace{-0.2cm}

Since
\[
\vect{u}\cdot\vect{v}
=
2(1)+1(-2)=0,
\]
the vectors are orthogonal.

\end{exampleblock}

\end{frame}


% =========================================================
% ORTHOGONAL PROJECTION
% =========================================================

\begin{frame}{Orthogonal Projection}

The projection of $\vect{u}$ onto a non-zero vector $\vect{v}$ is

\[
\boxed{
\operatorname{proj}_{\vect{v}}(\vect{u})
=
\frac{
\vect{u}\cdot\vect{v}
}{
\vect{v}\cdot\vect{v}
}
\vect{v}
}
\]

\vspace{0.5cm}

The projection gives the component of $\vect{u}$ that lies in the direction of $\vect{v}$.

\vspace{0.3cm}

It is useful in:

\begin{itemize}
    \item geometry
    \item least-squares methods
    \item computer graphics
    \item machine learning
\end{itemize}

\end{frame}


% =========================================================
% VECTOR DECOMPOSITION
% =========================================================

\begin{frame}{Decomposing a Vector}

A vector can be separated into components parallel and perpendicular to another vector.

\[
\boxed{
\vect{u}
=
\operatorname{proj}_{\vect{v}}(\vect{u})
+
\vect{u}_{\perp}
}
\]

where

\[
\vect{u}_{\perp}
=
\vect{u}
-
\operatorname{proj}_{\vect{v}}(\vect{u}).
\]

\vspace{0.4cm}

The two components satisfy

\[
\boxed{
\vect{u}_{\perp}\cdot\vect{v}=0.
}
\]

\end{frame}


% =========================================================
% DOT PRODUCT AND DISTANCE
% =========================================================

\begin{frame}{Dot Product and Distance}

The dot product can also be used to express the squared distance between two vectors.

For $\vect{u},\vect{v}\in\mathbb{R}^n$,

\[
\|\vect{u}-\vect{v}\|^2
=
(\vect{u}-\vect{v})^T
(\vect{u}-\vect{v}).
\]

Expanding,

\[
\boxed{
\|\vect{u}-\vect{v}\|^2
=
\vect{u}^T\vect{u}
+
\vect{v}^T\vect{v}
-
2\vect{u}^T\vect{v}
}
\]

Thus the dot product is closely connected to both:

\[
\text{length}
\qquad\text{and}\qquad
\text{distance}.
\]

\end{frame}


% =========================================================
% DOT PRODUCT AS A FUNCTION
% =========================================================

\begin{frame}{Dot Product and Duality}

For a fixed vector $\vect{w}$, consider

\[
f(\vect{x})
=
\vect{w}^{T}\vect{x}.
\]

This takes a vector $\vect{x}$ and produces a scalar.

\vspace{0.4cm}

\[
\boxed{
\vect{x}
\longrightarrow
\vect{w}^{T}\vect{x}
}
\]

Such a function is a \textbf{linear functional}.

\vspace{0.4cm}

This provides a connection between:

\[
\boxed{
\text{vectors}
\quad\longleftrightarrow\quad
\text{linear functions}
}
\]

and leads naturally to the idea of \textbf{duality}.

\end{frame}


% =========================================================
% WEIGHT VECTORS
% =========================================================

\begin{frame}{Weight Vectors}

A weight vector can define a linear function
\[
f(\vect{x})
=
\vect{w}^{T}\vect{x}.
\]

For
\[
\vect{w}
=
\begin{bmatrix}
w_1\\
w_2
\end{bmatrix},
\qquad
\vect{x}
=
\begin{bmatrix}
x_1\\
x_2
\end{bmatrix},
\]

we obtain
\[
f(\vect{x})
=
w_1x_1+w_2x_2.
\]

In machine learning, weight vectors are fundamental to:

\begin{itemize}
    \item linear models
    \item linear classifiers
    \item perceptrons
\end{itemize}

\end{frame}


% =========================================================
% HYPERPLANE
% =========================================================

\begin{frame}{Hyperplanes}

A vector $\vect{w}$ can define a hyperplane through

\[
\boxed{
\vect{w}^{T}\vect{x}=c
}
\]

where $c$ is a constant.

\vspace{0.2cm}

The vector $\vect{w}$ is \textbf{normal} to the hyperplane.

\vspace{0.2cm}

In two dimensions:

\[
w_1x_1+w_2x_2=c
\]

represents a line.

In three dimensions:
\[
w_1x_1+w_2x_2+w_3x_3=c
\]

represents a plane.

\end{frame}


% =========================================================
% CLASSIFICATION INTERPRETATION
% =========================================================

\begin{frame}{Hyperplanes in Classification}

A linear classifier can use

\[
f(\vect{x})
=
\vect{w}^{T}\vect{x}+b.
\]

The decision boundary is

\[
\boxed{
\vect{w}^{T}\vect{x}+b=0
}
\]

\vspace{0.4cm}

The space can then be divided according to the sign of $f(\vect{x})$:

\[
\begin{cases}
f(\vect{x})>0 & \text{one side}\\
f(\vect{x})<0 & \text{other side}.
\end{cases}
\]

\vspace{0.4cm}

This is the geometric foundation of models such as the \textbf{perceptron}.

\end{frame}